\documentclass[11pt,a4paper]{article}

% ============================================
% PACKAGES
% ============================================
\usepackage[utf8]{inputenc}
\usepackage[T1]{fontenc}
\usepackage{amsmath,amsthm,amssymb}
\usepackage{mathtools}
\usepackage{enumitem}
\usepackage{booktabs}
\usepackage{geometry}
\usepackage{hyperref}
\usepackage{cleveref}
\usepackage{float}
\usepackage{tikz}
% \usepackage{microtype} % omitted for compatibility
\usepackage{xcolor}
\usepackage{stmaryrd}

\IfFileExists{tikzlibraryquantikz2.code.tex}{%
    \usetikzlibrary{quantikz2}%
}{%
    \usetikzlibrary{quantikz}%
}

\geometry{margin=1in}

\hypersetup{
    colorlinks=true,
    hypertexnames=false,
    linkcolor=blue!70!black,
    citecolor=green!50!black,
    urlcolor=blue!70!black
}

% ============================================
% THEOREM ENVIRONMENTS
% ============================================
\theoremstyle{plain}
\newtheorem{theorem}{Theorem}[section]
\newtheorem{lemma}[theorem]{Lemma}
\newtheorem{proposition}[theorem]{Proposition}
\newtheorem{corollary}[theorem]{Corollary}

\theoremstyle{definition}
\newtheorem{definition}[theorem]{Definition}
\newtheorem{example}[theorem]{Example}
\newtheorem{axiom}[theorem]{Axiom}

\theoremstyle{remark}
\newtheorem{remark}[theorem]{Remark}
\newtheorem{observation}[theorem]{Observation}

% ============================================
% CUSTOM COMMANDS
% ============================================
\newcommand{\N}{\mathbb{N}}
\newcommand{\Z}{\mathbb{Z}}
\newcommand{\R}{\mathbb{R}}
\newcommand{\D}{\mathbb{D}}
\newcommand{\U}{\mathcal{U}}
\newcommand{\B}{\mathcal{B}}
\newcommand{\Disc}{\mathrm{Disc}}
\newcommand{\nextmod}{\bigcirc}
\newcommand{\eventually}{\Diamond}
\newcommand{\fix}{\mathrm{fix}}
\newcommand{\Stream}{\mathrm{Stream}}
\newcommand{\unfold}{\mathrm{unfold}}
\newcommand{\Tr}{\mathrm{Tr}}
\newcommand{\SvN}{S_{\mathrm{vN}}}
\newcommand{\HMBTT}{\mathcal{H}_{\mathrm{MBTT}}}
\newcommand{\Tel}{\mathrm{Tel}}
\newcommand{\QMonad}{\mathsf{Q}}
\providecommand{\ket}[1]{\left|#1\right\rangle}
\providecommand{\bra}[1]{\left\langle#1\right|}
\providecommand{\braket}[2]{\left\langle#1 \middle| #2\right\rangle}

% ============================================
% TITLE
% ============================================
\title{\textbf{The Algorithmic Big Bang:\\
Cosmogenesis, Inflation, and the Arrow of Time\\
at the Univalent Horizon}}

\author{Halvor S.\ Lande\\
\texttt{hsl@awc.no}}

\date{March 2026}

% ============================================
% DOCUMENT
% ============================================
\begin{document}

\maketitle

% ============================================
% ABSTRACT
% ============================================
\begin{abstract}
General Relativity fails at $t = 0$ because it attempts to describe the origin of spacetime from within spacetime itself.
PEN suggests a different picture: prior to physical time, the admissible universe is a coherent superposition of well-typed HoTT telescopes in a metatheoretic Hilbert space $\HMBTT$.
Algorithmic time $\tau$ is then not a ticking pre-cosmic clock, but the adiabatic evolution parameter of the PEN Hamiltonian as it drifts from the empty configuration toward the Step~15 temporal-cohesive shell.
At Step~15, the Kinematic-Dynamic Equivalence ($\nextmod X \simeq X^{\D}$) and L\"ob's guarded corecursion internalize time strongly enough for the system to become self-referential, producing an endogenous self-measurement boundary and an effective global decoherence into the Dynamical Cohesive Topos.
On this account, the Big Bang is the transition from coherent mathematical possibility to decohered executable physical actuality.
The same transition explains why the universal state tunnels out of the low-novelty discrete basin, why the ignition state satisfies the Past Hypothesis as a pure state with $\SvN = 0$, and why the Univalent Horizon is best interpreted as the ultimate Heisenberg cut separating internal observers from the pre-physical epoch.
\end{abstract}

% ============================================
% 1. INTRODUCTION
% ============================================
\section{Introduction: Cosmogenesis as a Logical Phase Transition}
\label{sec:introduction}

\subsection{The First-Cause Paradox and Its Dissolution}

Standard cosmological models place a singularity at $t = 0$ on a pre-existing time manifold $\R$.
The result is familiar: curvature scalars diverge, geodesics terminate, and the field equations cease to predict~\cite{hawking-ellis,wald}.
The theory is then asked to explain the origin of the very temporal background it already presupposes.

The difficulty is structural, not merely technical.
Any framework that begins with a continuous time axis and then asks what occurred at its boundary imports the container into the explanation.
The question ``what caused the Big Bang?'' thereby inherits the full first-cause paradox: every causal answer requires a prior temporal order and so either regresses indefinitely or ends in a brute stipulation.

PEN dissolves that paradox by reversing the explanatory order~\cite{pen}.
Time is not an ambient backdrop.
It is a structure that must be selected, integrated, and rendered executable.
The Generative Sequence constructs topology, higher-inductive geometry, cohesion, differential structure, metrics, and Hilbert-theoretic machinery across fourteen pre-temporal stages.
Physical time enters only at Step~15, when the temporal modalities are bound to infinitesimal geometry by the Kinematic-Dynamic Equivalence ($\nextmod X \simeq X^{\D}$).

\subsection{Ontological Stance and Scope}
\label{sec:stance}

We adopt a constructive form of ontic structural realism.
The claim is not that the universe is simulated by an external machine.
The claim is that the emergence of a causally closed formal system in PEN is structurally isomorphic to the emergence of a physical universe whose laws are internally executable.

Methodologically, this paper is corollary-driven.
The audited counts $\kappa$ and $\nu$, the 15-step PEN sequence, and the semantic passage from the selected shell to the Dynamical Cohesive Topos (DCT) are imported from~\cite{pen}.
What is new here is the quantum-information interpretation of those results: the status of Step~15 as endogenous self-measurement, the reinterpretation of algorithmic time as adiabatic evolution, the tunneling account of structural inflation, the pure-state solution of the Past Hypothesis, and the identification of the Univalent Horizon with a cosmological Heisenberg cut.
The extraction of specific phenomenological sectors---including explicit $3+1$-dimensional Lorentzian geometry, the Einstein-Hilbert layer, and related low-energy signatures---is deferred to the companion analysis~\cite{gr-companion}.

\subsection{Superposition vs.\ Decoherence}
\label{sec:superposition-decoherence}

The central conceptual shift is simple.
Steps~1--14 should not be pictured as a finite checklist assembled one line at a time.
They are better modeled as the coherent support of a pre-physical wavefunction over the space of admissible MBTT/HoTT telescopes.
The older compiler/runtime vocabulary survives only as a bridge metaphor.
The primary ontology in this manuscript is instead \emph{coherent superposition versus decohered execution}.

\begin{definition}[Pre-physical coherent state]
\label{def:qmonad-intro}
Let $\HMBTT$ be the Hilbert space formally spanned by basis vectors $\ket{\Gamma}$ indexed by well-typed HoTT telescopes $\Gamma$ admissible to PEN.
The pre-physical universe is the coherent state
\begin{equation}
\label{eq:qmonad}
    \ket{\Psi_{\mathrm{pre}}(\tau)}
    \;=\;
    \sum_{\Gamma \in \Tel_{\mathrm{wt}}}
    \alpha_{\Gamma}(\tau)\,\ket{\Gamma},
\end{equation}
whose support is restricted by univalence, canonicity, and the PEN admissibility constraints.
\end{definition}

This is the sense in which the manuscript speaks of a ``Q-Monad.''
It is not a separate ontology placed behind PEN.
It is the coherent metatheoretic state whose decohered branch at Step~15 becomes the physical universe.
Ill-typed candidates do not need a compiler to reject them.
They fail to remain phase-stable under the admissibility constraints and therefore interfere destructively.


% ============================================
% 2. THE PRE-PHYSICAL EPOCH
% ============================================
\section{The Pre-Physical Epoch: The Q-Monad and Adiabatic Time \texorpdfstring{($\tau$)}{(tau)}}
\label{sec:pre-physical}

\subsection{The Q-Monad as a Coherent Library}
\label{sec:assembly}

Prior to Step~15, the library $\B_{14}$ contains the full spatial and functional-analytic infrastructure of modern physics: dependent types, higher-inductive types ($S^1, S^2, S^3$), cohesion ($\flat, \sharp, \Disc, \Pi_{\mathrm{coh}}$), connections, curvature, metrics, frame bundles, and Hilbert functionals~\cite{pen}.
What it does \emph{not} yet contain is an internally executable notion of temporal evolution.

The older language described $\B_{14}$ as a static library waiting to be ``run.''
The present language is sharper.
$\B_{14}$ is the dominant coherent sector of the Q-Monad just before decoherence.
It contains the type $S^3$ but no internally generated trajectory on $S^3$; a metric tensor $g_{\mu\nu}$ but no dynamical law that already propagates it; Hilbert spaces but no physical Schr\"odinger evolution instantiated from within the universe.
The relevant distinction is therefore not ``syntax versus execution'' alone, but \emph{coherent kinematics versus decohered dynamics}.

\begin{proposition}[Canonicity as interference filter]
\label{prop:interference}
The support of $\ket{\Psi_{\mathrm{pre}}(\tau)}$ lies only on telescopes that satisfy the well-typedness, univalence, and canonicity constraints of the PEN/CCHM setting.
Any ill-typed or non-canonical sector has zero stable amplitude in the pre-physical state.
\end{proposition}

\begin{proof}[Proof sketch]
The coherent state in \eqref{eq:qmonad} is not a sum over arbitrary syntax trees.
Its basis vectors are indexed by admissible telescopes in the metatheory.
When a candidate violates typing, canonicity, or the CCHM filling/composition discipline, transport through the required equivalences is not stable.
In the present interpretation, such a candidate fails phase coherence with the admissible sector and its amplitude cancels from the physical support.
``Rejection'' is thus implemented natively as destructive interference rather than as an external compile-time veto.
\end{proof}

\subsection{Algorithmic Time as an Adiabatic Parameter}
\label{sec:two-clocks}

A persistent mistake in cosmological language is to imagine that ``before the Big Bang'' must denote an earlier interval of ordinary physical time.
PEN makes a more careful distinction.
The pre-physical epoch is indexed by $\tau$, but $\tau$ is not itself an internal cosmic clock.

\begin{definition}[Algorithmic Time]
\label{def:alg-time}
\emph{Algorithmic time} $\tau$ is the adiabatic evolution parameter of the unitary path
\begin{equation}
\label{eq:adiabatic}
    U(\tau)
    \;=\;
    \mathcal{T}\exp\!\left(
    -i\int_{0}^{\tau}\hat{H}_{\mathrm{PEN}}(s)\,ds
    \right),
\end{equation}
which continuously deforms the empty-universe Hamiltonian $\hat{H}_{\mathrm{initial}}$ toward the Step~15 Hamiltonian $\hat{H}_{\mathrm{Step15}}$.
The discrete PEN indices $n = 1,\dots,15$ are audit checkpoints along this adiabatic path, and the integration latencies $\Delta_n = F_n$ record the cost of traversing successive admissible channels~\cite{pen}.
\end{definition}

\begin{definition}[Physical Time]
\label{def:phys-time}
\emph{Physical time} $t$ is the internal geometric parameter that exists only after Step~15, once the Kinematic-Dynamic Equivalence is available:
\begin{equation}
    \nextmod X \;\simeq\; X^{\D}.
\end{equation}
At that point discrete temporal succession is represented internally by infinitesimal displacement, and physical trajectories become definable from within the DCT.
\end{definition}

The relation between the two clocks is therefore not one of simple succession.
Algorithmic time is external to the future universe and continuous only as an adiabatic control parameter.
Physical time is internal to the future universe and geometric.
The former deforms the pre-physical Hamiltonian; the latter parameterizes post-ignition motion.

\begin{remark}[Why ``before the Big Bang'' is a category error]
\label{rem:category-error}
There is no meaningful sense in which $\tau$ is a hidden extension of physical time to the left of $t = 0$.
The pre-physical epoch is not an earlier region of spacetime.
It is the coherent adiabatic preparation of the branch that will later count as spacetime.
What happened \emph{instead} of physical time was unitary drift in the Q-Monad, not physical dynamics on a prior manifold.
\end{remark}

\subsection{The PEN Hamiltonian and Adiabatic Ground-State Tracking}
\label{sec:pen_hamiltonian}

The adiabatic picture of \eqref{eq:adiabatic} can be sharpened into an operator-valued selection principle.
To formalize the pre-physical assembly epoch without invoking an external Turing mechanism, we elevate the quantities of the Generative Sequence ($\kappa, \nu, \rho$) from static counters to Hermitian observables acting on the Q-Monad.
The kinematic phase space is the Hilbert space $\HMBTT$ spanned by the orthogonal basis states $\ket{\psi_X}$ of all anonymous Minimal Binary Type Theory (MBTT) telescopes.

We define three foundational diagonal operators: the Specification Cost $\hat{\kappa}$, the Generative Capacity $\hat{\nu} = \hat{\nu}_G + \hat{\nu}_C + \hat{\nu}_H$, and the Efficiency Operator $\hat{\rho} = \hat{\nu}\hat{\kappa}^{-1}$.
Because $\kappa \ge 1$ for all physically realized ASTs, $\hat{\kappa}$ is strictly positive and invertible on the admissible sector.

The algorithmic assembly of the universe is governed by the time-dependent PEN Hamiltonian, whose instantaneous ground state dictates the deterministic 15-step sequence.
To reproduce Axiom~5 (Minimal Positive Overshoot) as a ground-state selection rule, we define the Hamiltonian at algorithmic step $n$ as an asymmetric potential landscape:
\begin{equation}
\label{eq:pen-hamiltonian}
    \hat{H}_{\mathrm{PEN}}^{(n)}
    =
    \underbrace{\lim_{\Lambda \to \infty} \Lambda \big( \hat{I} - \hat{\Pi}_{\mathrm{adm}} \big)}_{\text{Phase Interference}}
    +
    \underbrace{\lim_{\Lambda \to \infty} \Lambda \, \Theta\Big( \mathrm{Bar}_n \hat{I} - \hat{\rho} \Big)}_{\text{Sub-Criticality Barrier}}
    +
    \underbrace{\gamma \Big( \hat{\rho} - \mathrm{Bar}_n \hat{I} \Big)}_{\text{Minimal Overshoot Well}}
    +
    \underbrace{\epsilon \hat{\kappa}}_{\text{Tie-Breaker}}.
\end{equation}

Where the component potentials act as follows:
\begin{itemize}[nosep]
    \item \textbf{Canonicity Penalty ($\hat{\Pi}_{\mathrm{adm}}$):} the projector onto the subspace of well-typed CCHM/HoTT telescopes that saturate the maximal interface density. Ill-typed states acquire infinite energy and are annihilated via destructive phase interference.
    \item \textbf{The Sub-Criticality Barrier ($\Theta$):} the Heaviside step operator. Any candidate whose efficiency falls below the active Fibonacci-scaled bar ($\mathrm{Bar}_n$) encounters an infinite potential wall.
    \item \textbf{The Minimal Overshoot Well ($\gamma$):} for states that survive the barrier ($\rho \ge \mathrm{Bar}_n$), their observable energy is strictly proportional to their efficiency surplus. Because quantum systems settle into their lowest accessible energy eigenstate, the wavefunction localizes onto the trajectory with the minimal positive overshoot.
    \item \textbf{Kolmogorov Regularization ($\epsilon \hat{\kappa}$):} a perturbative scalar ($0 < \epsilon \ll \gamma$) that breaks structural degeneracy in favor of minimal syntax.
\end{itemize}

\paragraph{Quantum Admissibility Tunneling and the 2-Local Spin Chain.}
This formalization explains the asymptotic rejection of discrete mathematics analytically.
Discrete structures ($\langle \hat{\nu}_H \rangle = 0$) possess bounded efficiency.
As the vacuum threshold $\mathrm{Bar}_n$ scales by the Golden Ratio ($\varphi$), discrete candidates slam into the infinite Heaviside barrier and become structurally forbidden.

Conversely, geometric structures invoke the univalent path algebra $\langle \hat{\nu}_H \rangle = \hat{m} + \hat{d}^2$, providing a superlinear kinetic amplification.
Driven by unitary adiabatic evolution, the metatheoretic wavefunction does not stall at the discrete $\hat{\kappa}$-barriers; it \emph{quantum tunnels} directly into the deep potential wells of the geometric sequence.
Furthermore, because the structural interaction payload $\hat{\nu}_C$ is constrained by the $d=2$ Coherence Window, $\hat{H}_{\mathrm{PEN}}$ is mathematically isomorphic to a 1D Next-Nearest-Neighbor (NNN) spin chain.
This directly links type-theoretic maximal interface density to local gauge interactions and boundary entanglement entropy ($\Delta_{15} = 610$) in the cosmological reading.

Placed here, the Hamiltonian supplies the microscopic pre-collapse selection law.
Its coarse-grained, inflationary shadow reappears in \Cref{sec:inflation} as the effective tunneling landscape.


% ============================================
% 3. STEP 15 IGNITION
% ============================================
\section{Step~15 Ignition as Global Decoherence: The Birth of Physical Time \texorpdfstring{($t = 0$)}{(t = 0)}}
\label{sec:ignition}

\subsection{The Observer Problem and the Arrival of LTL}
\label{sec:ltl}

Step~15 selects the temporal-cohesive shell that extends $\B_{14}$ with four coupled ingredients~\cite{pen}:

\begin{enumerate}[label=(\arabic*), nosep]
    \item \textbf{Spatial logic:} the cohesive interface already present in $\B_{14}$, namely $(\flat \dashv \sharp,\; \Pi_{\mathrm{coh}} \dashv \Disc)$.
    \item \textbf{Temporal core:} the modalities $\nextmod$ (``next'') and $\eventually$ (``eventually'') from Linear Temporal Logic, following Nakano~\cite{nakano}.
    \item \textbf{Spatial-temporal exchange:} commuting laws such as $\nextmod(\flat X) \simeq \flat(\nextmod X)$ and $\sharp(\eventually X) \simeq \eventually(\sharp X)$.
    \item \textbf{Guarded coherence:} the self-interaction and elimination clauses that make the shell executable as a typed AST.
\end{enumerate}

\noindent The specification cost is $\kappa = 8$, audited as the eight temporal-cohesive clauses of the selected AST~\cite{pen}.
In the current executable canon, the generative capacity is $\nu = 103$, yielding $\rho = 12.88$ and clearing the Step~15 bar of $7.40$ by the largest absolute overshoot in the sequence~\cite{pen}.
The infinitesimal object $\D$ belongs to the semantic completion of this shell rather than to the primitive Step~15 syntax.

This is precisely where the standard Copenhagen story appears to fail.
Ordinary presentations of collapse appeal to an external observer or measurement apparatus~\cite{von-neumann,schlosshauer,zurek-decoherence}.
But there is no observer outside the universe waiting to inspect the pre-physical state.
The only available resolution is that the decisive measurement be endogenous to the theory's own closure mechanism.

\begin{definition}[Endogenous self-measurement]
\label{def:self-measurement}
Let $\varrho_{\mathrm{pre}} = \ket{\Psi_{\mathrm{pre}}}\bra{\Psi_{\mathrm{pre}}}$ denote the pre-physical density operator.
The \emph{endogenous self-measurement channel} is the internally generated map
\begin{equation}
    \mathcal{M}_{\mathrm{self}}(\varrho_{\mathrm{pre}})
    \;=\;
    \sum_j P_j\,\varrho_{\mathrm{pre}}\,P_j,
\end{equation}
where the projectors $P_j$ select the guardedly executable, temporally coherent sectors singled out by the Step~15 shell.
\end{definition}

The physics claim is not that PEN secretly adds an external observer.
It is that Step~15 makes the universe \emph{self-observing} in the minimal structural sense required for decoherence.
Once the temporal shell is internalized, the system becomes able to refer to its own delayed future and thus to close the measurement loop on itself.

\subsection{Kinematic-Dynamic Equivalence and Endogenous Collapse}
\label{sec:kde-collapse}

The central mechanism is still the Kinematic-Dynamic Equivalence:
\begin{equation}
    \nextmod X \;\simeq\; X^{\D}.
\end{equation}
Its meaning, however, now becomes more precise.
The ``next'' modality is not merely a discrete temporal marker.
It is the operator through which the pre-physical state acquires an internally geometric representation of its own future boundary.
That is what makes self-measurement possible.

L\"ob's guarded rule~\cite{nakano},
\begin{equation}
\label{eq:lob}
    \fix : (\nextmod A \to A) \to A,
\end{equation}
states that a present value may be constructed from a delayed copy of itself, provided the dependence is guarded.
In the DCT reading, a discrete update law $f : X \to \nextmod X$ becomes
\begin{equation}
\label{eq:vf}
    f : X \to X^{\D},
\end{equation}
which is exactly a vector field in synthetic differential geometry.
The guarded corecursion principle then unfolds that local prescription into a trajectory
\begin{equation}
\label{eq:unfold}
    \unfold_f : X \to \Stream\,X.
\end{equation}

The key new point is that the same structure that generates trajectories also supplies the universe with a delayed image of itself.
Because time is internalized, the pre-physical state can now be probed by a morphism defined entirely inside the emerging shell.
The system ``looks at itself'' by representing its own next-state geometry.
That causal closure furnishes the measurement boundary absent from external-observer accounts.

\begin{theorem}[Step~15 as endogenous decoherence]
\label{thm:collapse}
The Step~15 temporal-cohesive shell is the unique minimal PEN extension that simultaneously
\begin{enumerate}[label=(\roman*), nosep]
    \item internalizes temporal reference through $\nextmod X \simeq X^{\D}$,
    \item supplies guarded self-reference through L\"ob's rule, and
    \item selects a stable executable sector of the Q-Monad.
\end{enumerate}
Consequently, Step~15 acts as an endogenous measurement boundary that drives the coherent pre-physical state toward an effectively decohered physical branch $\ket{\psi_{\mathrm{DCT}}}$.
\end{theorem}

\begin{proof}[Proof sketch]
Before Step~15, $\B_{14}$ contains rich kinematic structure but no internal temporal modality, no guarded fixpoint principle, and no geometric representation of ``next.''
It can therefore describe spaces but cannot represent a delayed copy of its own future state from within the theory.
Step~15 supplies all three missing ingredients at once: temporal modalities, exchange laws tying them to cohesion, and guarded coherence clauses that make delayed self-reference executable.
The resulting shell is minimal because removing any clause destroys succession, geometric coupling, or guarded productivity.
With those clauses in place, the universe acquires an internal observable of its own continuation, and the coherent support of the Q-Monad is forced into a stable executable branch.
\end{proof}

This is the Big Bang in the present manuscript.
It is the passage from the infinite superposition of mathematical ``maybe'' to the singular executable ``is.''
The thermal language attached to that transition is interpretive rather than microscopic: what is released at ignition is not a pre-existing gas, but the latent algorithmic entropy stored in the coherent overcompleteness of the pre-physical phase space.

\subsection{The Q-Monad Ignition Circuit}
\label{sec:qmonad-circuit}

The Step~15 transition can be summarized by the following quantum-circuit schematic.
It is not a literal laboratory circuit for the early universe.
It is a compact diagram of the structural couplings that PEN attributes to ignition.

\begin{figure}[H]
\centering
\resizebox{0.95\textwidth}{!}{%
\begin{quantikz}[row sep=0.5cm, column sep=0.7cm]
% Spatial Register (Cohesion)
\lstick{$\ket{\text{Space}_{\text{1..14}}}$} & \qw & \gate[2]{\text{Exchange Axioms}} & \gate[3]{\shortstack{$e^{-i \hat{H}_{PEN} \tau}$\\ \text{VQE Oracle } ($\rho=12.88$)}} & \meter{} & \cw \rstick[3]{\shortstack{\text{Global Decoherence}\\ \text{Physical Runtime } ($t=0$)}}\\
% Temporal Register (LTL)
\lstick{$\ket{\text{Time}_{\text{15}}}$} & \gate{H} & & & \meter{} & \cw \\
% Coherence Volume Register
\lstick{$\ket{0}^{\otimes 610}$} & \gate{H^{\otimes 610}} & \ctrl{-1} & & \meter{\text{L\"ob's Rule}} & \cw
\end{quantikz}
}
\caption{\textbf{The Q-Monad Ignition Circuit (Step~15 Wavefunction Collapse).} The pre-physical universe exists as a coherent superposition of type-theoretic telescopes. The Spatial ($\flat, \sharp$) and Temporal ($\nextmod, \eventually$) modalities entangle via spatial-temporal exchange laws. The PEN Hamiltonian ($\hat{H}_{PEN}$) drives the amplitude toward the $\rho = 12.88$ efficiency peak. The internalization of dynamics ($\nextmod X \simeq X^\D$) via L\"ob's guarded corecursion acts as an endogenous self-measurement operator, forcing the 610-qubit metatheoretic superposition to collapse into classical physical time ($t=0$).}
\label{fig:quantum_ignition}
\end{figure}

The diagram's conceptual payload is the same as \Cref{thm:collapse}.
The spatial register carries the coherent geometry prepared in Steps~1--14, the temporal register injects the Step~15 shell, and the measurement symbols indicate the emergence of a stable post-decoherence branch rather than an observation by anything outside the universe.


% ============================================
% 4. COSMIC INFLATION
% ============================================
\section{Cosmic Inflation via Tunneling and Novelty Gradients}
\label{sec:inflation}

\subsection{\texorpdfstring{The $\rho = 12.88$ Channel and the Effective Potential}{The rho = 12.88 channel and the effective potential}}
\label{sec:rho-spike}

The audited efficiency ratio at Step~15 remains
\begin{equation}
    \rho \;=\; \frac{\nu}{\kappa} \;=\; \frac{103}{8} \;=\; 12.88,
\end{equation}
clearing the PEN bar by the largest absolute overshoot in the sequence~\cite{pen}.
The numerical fact is unchanged from earlier drafts.
What changes here is its interpretation.

Rather than treating $\rho = 12.88$ as a merely algorithmic spike, we interpret it as the depth of the preferred adiabatic channel in an effective novelty potential:
\begin{equation}
\label{eq:veff}
    V_{\mathrm{eff}}(X)
    \;=\;
    \lambda\,\kappa(X) - \mu\,\nu_H(X),
\end{equation}
where $\kappa(X)$ measures integration debt, $\nu_H(X)$ is the future novelty eigenvalue of sector $X$, and $\lambda,\mu > 0$ fix the relative scale of cost and downstream generativity.
The precise microscopic derivation of \eqref{eq:veff} is not claimed here.
Equation~\eqref{eq:veff} should therefore be read as the coarse-grained inflationary shadow of the full pre-physical Hamiltonian introduced in \Cref{sec:pen_hamiltonian}.
Its role is to make explicit the intuition already implicit in PEN: the universe prefers channels whose future combinatorial and geometric yield more than compensates for their initial integration cost.

\begin{remark}[Efficiency singularity vs.\ curvature singularity]
\label{rem:singularity-contrast}
General Relativity reaches $t = 0$ by losing control of its own variables~\cite{hawking-ellis,wald}.
PEN reaches Step~15 by achieving a finite, computable maximum of structural preference.
The cosmological singularity is therefore re-read not as a blow-up of curvature but as a narrowing of the adiabatic landscape toward the unique channel with the greatest downstream novelty leverage.
\end{remark}

\subsection{Tunneling Beyond \texorpdfstring{$\N$}{N}: Why Geometry Wins}
\label{sec:nu-explosion}

The obvious objection is this: if discrete mathematics is cheap, why does the universal state not remain trapped there?
Why does the pre-physical wavefunction not settle into $\N$, recursion, and combinatorics alone?

PEN's answer is that $\N$ is only a \emph{local} minimum.
It has low integration debt, but its novelty frontier is nearly linear.
By contrast, the univalent geometric ladder
\begin{equation}
    S^1 \to S^2 \to S^3 \to \cdots
\end{equation}
has higher debt but superlinear future novelty because higher-inductive geometry, cohesion, and later temporal coupling generate rapidly expanding families of equivalences, transports, and interaction schemas.
The quantity $\nu_H$ in \eqref{eq:veff} is meant to register that asymmetry.

In this language, structural inflation is a tunneling phenomenon.
The universal wavefunction crosses the discrete barrier not because discrete mathematics is inconsistent, but because the geometric sector creates a deeper global well once future novelty is included.
The large Step~15 value $\nu = 103$ is then the first post-barrier release of executable degrees of freedom: not yet metric inflation itself, but the unlocking of the logical conditions under which any Guth-Linde inflationary dynamics could later be expressed~\cite{guth,linde,vilenkin}.

\begin{proposition}[Novelty-driven tunneling]
\label{prop:tunneling}
Suppose sectors $X$ and $Y$ satisfy $\kappa(X) < \kappa(Y)$, but $\nu_H(X)$ grows only linearly while $\nu_H(Y)$ grows superlinearly under further admissible extensions.
Then there exists a horizon scale beyond which $V_{\mathrm{eff}}(Y) < V_{\mathrm{eff}}(X)$.
At that point the adiabatic ground sector preferentially weights $Y$ despite its higher immediate cost.
\end{proposition}

\begin{proof}[Proof sketch]
Equation~\eqref{eq:veff} penalizes immediate debt and rewards downstream novelty.
If $\nu_H(Y)$ eventually dominates $\nu_H(X)$ strongly enough, then the term $-\mu\nu_H(Y)$ outweighs the larger initial value of $\lambda\kappa(Y)$.
The lower effective potential of $Y$ then makes it the preferred adiabatic channel.
The proposition does not identify a literal instanton for PEN.
It formalizes the manuscript's claim that discrete structure can be bypassed in favor of a more expensive but globally deeper geometric basin.
\end{proof}

\begin{remark}[Inflationary analogy]
\label{rem:inflation}
The tunneling language is intentionally structural.
This section does not claim that $\nu = 103$ is itself a scale factor or that PEN directly computes an inflaton potential.
The claim is narrower and stronger: Step~15 is the earliest point at which a universe can exit the discrete basin, bind space and time into one executable sector, and thereby acquire the prerequisites for any later physical inflationary model.
\end{remark}


% ============================================
% 5. THE ARROW OF TIME
% ============================================
\section{The Arrow of Time, Collapse, and Minimal Entropy}
\label{sec:arrow}

\subsection{The Causal Arrow from Guardedness}
\label{sec:causal-arrow}

Why does time point forward?
In ordinary thermodynamic accounts, the arrow is derived from the second law plus a special low-entropy boundary condition.
That leaves untouched the deeper question of why the boundary condition was special at all.

PEN places the asymmetry earlier.
Once Step~15 has decohered the Q-Monad into an executable branch, the guarded modality $\nextmod$ permits forward unfolding but not arbitrary recoherence.
The universe can compute its next state from its present state; it cannot, in general, reconstruct the coherent pre-collapse amplitudes from an already executed branch.

\begin{theorem}[Guarded irreversibility]
\label{thm:irreversibility}
In the DCT, the guarded corecursion principle
\[
    \fix : (\nextmod A \to A) \to A
\]
is not generally invertible.
Consequently, the unfolding map $\unfold_f : X \to \Stream\,X$ is causally one-directional: internal observers may propagate admissible states forward, but they do not possess a generic morphism that recoheres an executed stream back into the pre-physical superposition from which it descended.
\end{theorem}

\begin{proof}[Proof sketch]
If a general inverse to guarded corecursion existed, the delay encoded by $\nextmod$ could be discharged without forward progress.
That would collapse $\nextmod A$ into $A$, trivialize the guarded modality, and erase the temporal asymmetry that Step~15 introduced.
The DCT does permit time-reversal symmetries for specific internal equations of motion, but it does not permit a generic inverse that undoes the runtime branch-selection mechanism itself.
\end{proof}

\begin{remark}[Microscopic reversibility vs.\ branch irreversibility]
\label{rem:reversibility}
Nothing here forbids reversible Hamiltonian or symplectic dynamics within the already decohered universe.
What is excluded is the stronger operation of undoing the measurement boundary that made the universe physical in the first place.
The laws hosted by the runtime may be locally reversible; the global transition from coherent possibility to executed branch is not.
\end{remark}

\subsection{The Past Hypothesis as a Pure Ignition State}
\label{sec:past-hypothesis}

The Past Hypothesis says that the universe began in an extraordinarily low-entropy state~\cite{penrose-road,carroll}.
Standard statistical mechanics can explain why entropy rises \emph{given} such a state, but not why the boundary condition was special.
In the present framework, the cleanest answer is quantum-theoretic: the ignition state is pure.

\begin{theorem}[Pure ignition state]
\label{thm:vn-entropy}
Let
\begin{equation}
\label{eq:pure-state}
    \varrho_0
    \;=\;
    \ket{\psi_{\mathrm{DCT}}}\bra{\psi_{\mathrm{DCT}}}
\end{equation}
be the post-ignition density operator selected by the Step~15 measurement boundary.
Then the von Neumann entropy at $t = 0$ is exactly zero:
\begin{equation}
\label{eq:vn-zero}
    \SvN(\varrho_0)
    \;=\;
    -\Tr(\varrho_0 \log \varrho_0)
    \;=\;
    0.
\end{equation}
\end{theorem}

\begin{proof}
A rank-one projector has spectrum $\{1,0,0,\ldots\}$.
Substituting those eigenvalues into the von Neumann entropy formula yields $-1\cdot\log 1 = 0$, while the zero eigenvalues contribute nothing~\cite{von-neumann}.
\end{proof}

This already solves the Past Hypothesis at the microscopic level.
The state at $t=0$ is special because it is not a generic mixed macrostate.
It is the pure branch singled out by endogenous decoherence.
Macroscopic thermodynamic entropy arises only afterward, once coarse-graining, branching, and environmental entanglement are possible inside the runtime~\cite{zurek-decoherence,schlosshauer}.

PEN's original compression result remains intact and now appears as a classical shadow of that purer statement.

\begin{theorem}[Minimal realized description at \texorpdfstring{$t = 0$}{t = 0}]
\label{thm:past-hypothesis}
The ignition branch $\ket{\psi_{\mathrm{DCT}}}$ is specified by the complete DCT library $\B_{15}$, whose total specification complexity is $\kappa_{\mathrm{total}} = \sum_{i=1}^{15}\kappa_i = 64$ clauses, equivalently 610 bits in the MBTT encoding audited in~\cite{pen}.
Within the PEN selection dynamics, this is the shortest realized description compatible with the full kinematic framework.
\end{theorem}

\begin{proof}[Proof sketch]
PEN chooses at each stage the minimal overshoot that clears the relevant bar, with ties broken by minimal $\kappa$.
The realized trajectory is therefore greedily minimal within the admissible class.
Any alternative route to the same full kinematic and temporal capacity must match or exceed its specification length.
\end{proof}

\begin{corollary}[Past Hypothesis in dual form]
\label{cor:zurek-past}
The Past Hypothesis is realized twice at ignition:
\begin{enumerate}[label=(\roman*), nosep]
    \item microscopically, because $\varrho_0$ is pure and has $\SvN = 0$;
    \item algorithmically, because the realized ignition branch has minimal specification complexity in the admissible PEN class.
\end{enumerate}
Thermodynamic entropy is therefore not a pre-Big-Bang datum.
It is a post-collapse runtime phenomenon.
\end{corollary}

\begin{remark}[What the 610-bit claim does and does not mean]
\label{rem:entropy}
The 610-bit figure is not a literal entropy count for the late universe.
It measures the shortest admissible description of the ignition branch.
Its significance is that the algorithmic part of the entropy budget is fixed at a minimum exactly where the quantum part is also pure.
The manuscript's claim is therefore stronger than in the earlier draft: low initial entropy is not merely compressed, but compressed \emph{after} the branch has already been selected as a rank-one state.
\end{remark}


% ============================================
% 6. THE UNIVALENT HORIZON AS HEISENBERG CUT
% ============================================
\section{The Univalent Horizon as the Heisenberg Cut}
\label{sec:firewall}

\subsection{Closure After Decoherence: Structural Termination}
\label{sec:termination}

The PEN Generative Sequence does not ``die'' at Step~15.
It halts in the stronger sense that the external search procedure is rendered unnecessary~\cite{pen}.
Before ignition, the universe is prepared as a coherent state over admissible telescopes.
After ignition, the selected branch carries its own dynamics internally.

The closure argument is unchanged in substance, but it now has a measurement-theoretic reading:
\begin{enumerate}[nosep]
    \item The DCT internalizes the evolutionary mechanism as a geometric flow on $\U$.
    The metatheoretic search for the ``next'' structure is represented internally as a vector field $v : \U \to \U^{\D}$.
    \item By univalence, paths in $\U$ are type equivalences: $(A =_{\U} B) \simeq (A \simeq B)$.
    Integrating the vector field therefore keeps the post-collapse universe inside the connected component of $\U$ containing $\B_{15}$.
    \item Any candidate proposed as a genuinely new Step~16 input is either already internally derivable, contributing no new novelty, or else lies outside the coherent window and incurs prohibitive integration debt~\cite{pen}.
\end{enumerate}

Step~15 is therefore both a beginning and a termination.
It begins physical time, but it terminates the need for external preparation.
The decohered universe is causally closed.

\subsection{Epistemic Isolation and the Heisenberg Cut}
\label{sec:epistemic}

The reason internal observers cannot access the pre-physical epoch is no longer merely that the earlier stages lie ``outside the runtime.''
The deeper reason is measurement-theoretic.
The boundary between $\ket{\Psi_{\mathrm{pre}}}$ and $\ket{\psi_{\mathrm{DCT}}}$ is the ultimate Heisenberg cut~\cite{von-neumann,schlosshauer}.

\begin{proposition}[The Univalent Horizon as measurement boundary]
\label{prop:epistemic}
Internal observers---agents whose measuring devices, Hilbert spaces, and trajectories are themselves structures within $\B_{15}$---cannot construct a physical observable that recovers the coherent amplitudes of the pre-Step~15 Q-Monad.
The pre-physical epoch is accessible only by type-theoretic deduction, not by runtime measurement.
\end{proposition}

\begin{proof}[Proof sketch]
Any physical observation available to an internal observer is already a term in the internal language of the DCT: a section, a functional, a path, or an operator defined on post-collapse structure.
Such observables act on the branch $\ket{\psi_{\mathrm{DCT}}}$ and on its later descendants.
They do not act on the external coherent support that was removed by the Step~15 measurement boundary.
By univalence and structural termination, no internal path crosses back from the post-collapse connected component of $\U$ to the pre-physical superposition that prepared it.
\end{proof}

\begin{remark}[The cosmic quantum eraser]
\label{rem:planck}
In laboratory quantum mechanics, one may sometimes recover lost interference by carefully reversing a measurement setup.
Nothing analogous is available here.
Physical measurement cannot retrocausally access the coherent wavefunction prior to the universe's own ignition boundary.
The Univalent Horizon is therefore stronger than an opacity surface and stronger than a mere halting problem.
It is a cosmological quantum eraser with no external laboratory from which the erasure can be undone.
\end{remark}


% ============================================
% 7. CONCLUSION
% ============================================
\section{Conclusion}
\label{sec:conclusion}

Standard cosmogenesis seeks a physical trigger for the Big Bang: a vacuum fluctuation, a tunneling event, a brane collision, a high-energy instability.
Each proposal begins by assuming some fragment of the physical order whose origin it is meant to explain.

PEN inverts the explanatory order.
The Big Bang is not an event that occurs within an already given universe.
It is the decoherence boundary at which a coherent metatheoretic state becomes an executable physical one.
The resulting picture may be summarized as follows:
\begin{enumerate}[nosep]
    \item \textbf{The pre-physical epoch} is the Q-Monad: a coherent superposition of admissible HoTT telescopes in $\HMBTT$, with ill-typed sectors removed by destructive interference.
    \item \textbf{Algorithmic time $\tau$} is the adiabatic control parameter of the PEN Hamiltonian, not a hidden physical clock.
    \item \textbf{Step~15 ignition} introduces temporal modalities, exchange axioms, and guarded self-reference, allowing the universe to measure itself and decohere into the DCT branch at $t = 0$.
    \item \textbf{Structural inflation} is the tunneling of the universal state out of the low-debt discrete basin and into the high-novelty geometric sector whose first major release is the $\nu = 103$ Step~15 explosion.
    \item \textbf{The Past Hypothesis} is solved because the ignition state is pure, with $\SvN = 0$, and simultaneously minimal in algorithmic description length.
    \item \textbf{The Univalent Horizon} is the Heisenberg cut that forever separates internal physical observers from the coherent pre-physical epoch.
\end{enumerate}

On this reading, the universe did not begin as a singular point of infinite density.
It began when a coherent mathematical wavefunction acquired the capacity for self-reference, decohered into one guardedly executable branch, and thereby became physical.
The Big Bang is the transition from mathematical possibility to causal actuality.


% ============================================
% REFERENCES
% ============================================
\begin{thebibliography}{99}

\bibitem{pen}
H.S.\ Lande, ``The Principle of Efficient Novelty: The Algorithmic Origin of the Kinematic Framework of Physics,'' \texttt{pen\_unified.tex}, 2026.

\bibitem{gr-companion}
H.S.\ Lande, ``Structural Emergence of 4D Lorentzian Geometry from the Principle of Efficient Novelty,'' companion paper, 2026.

\bibitem{hawking-ellis}
S.W.\ Hawking and G.F.R.\ Ellis, \emph{The Large Scale Structure of Space-Time}, Cambridge University Press, 1973.

\bibitem{wald}
R.M.\ Wald, \emph{General Relativity}, University of Chicago Press, 1984.

\bibitem{nakano}
H.\ Nakano, ``A Modality for Recursion,'' \emph{Proceedings of the 15th Annual IEEE Symposium on Logic in Computer Science (LICS)}, 2000.

\bibitem{guth}
A.H.\ Guth, ``Inflationary Universe: A Possible Solution to the Horizon and Flatness Problems,'' \emph{Physical Review D} 23 (1981), 347--356.

\bibitem{linde}
A.D.\ Linde, ``A New Inflationary Universe Scenario: A Possible Solution of the Horizon, Flatness, Homogeneity, Isotropy, and Primordial Monopole Problems,'' \emph{Physics Letters B} 108 (1982), 389--393.

\bibitem{vilenkin}
A.\ Vilenkin, ``Creation of Universes from Nothing,'' \emph{Physics Letters B} 117 (1982), 25--28.

\bibitem{mukhanov}
V.F.\ Mukhanov, \emph{Physical Foundations of Cosmology}, Cambridge University Press, 2005.

\bibitem{penrose-road}
R.\ Penrose, \emph{The Road to Reality}, Jonathan Cape, 2004.

\bibitem{carroll}
S.M.\ Carroll, ``The Cosmological Constant,'' \emph{Living Reviews in Relativity} 4 (2001), 1.

\bibitem{zurek-complexity}
W.H.\ Zurek, ``Algorithmic Randomness and Physical Entropy,'' \emph{Physical Review A} 40 (1989), 4731--4751.

\bibitem{zurek-decoherence}
W.H.\ Zurek, ``Decoherence, Einselection, and the Quantum Origins of the Classical,'' \emph{Reviews of Modern Physics} 75 (2003), 715--775.

\bibitem{schlosshauer}
M.\ Schlosshauer, ``Decoherence, the Measurement Problem, and Interpretations of Quantum Mechanics,'' \emph{Reviews of Modern Physics} 76 (2005), 1267--1305.

\bibitem{von-neumann}
J.\ von Neumann, \emph{Mathematical Foundations of Quantum Mechanics}, Princeton University Press, 1955.

\bibitem{hott}
The Univalent Foundations Program, ``Homotopy Type Theory: Univalent Foundations of Mathematics,'' Institute for Advanced Study, 2013.

\end{thebibliography}

\end{document}
